\documentclass[11pt]{article}
\usepackage{amsmath,amssymb,amsthm,amsfonts}
\usepackage{geometry}
\geometry{margin=1in}
\usepackage{algorithm}
\usepackage{algpseudocode}

\newtheorem{theorem}{Theorem}
\newtheorem{lemma}{Lemma}
\newtheorem{assumption}{Assumption}
\newtheorem{corollary}{Corollary}
\newtheorem{definition}{Definition}

\title{Recursive Variance-Reduced Gradient Method with Norm-Adaptive Step Size\\
\large (A SARAH/SPIDER-type Algorithm)}
\author{}
\date{}

\begin{document}
\maketitle

%==========================================================
\section*{Algorithm Name}
We analyze a \emph{Recursive Variance Estimator with Norm-Adaptive Step Size} (RVE-NA),
a member of the SARAH/SPIDER family of variance-reduced stochastic gradient methods for
non-convex smooth optimization. The algorithm maintains an exponentially-weighted recursive
gradient estimator
\[
v_t = (1-\alpha)\, v_{t-1} + \alpha\, \nabla f(x_t),
\]
employs a step size that depends on the norm of $v_t$, and applies bias correction to
$v_t$ analogous to Adam's moment correction.

%==========================================================
\section*{Mathematical Formulation}

We aim to solve
\[
\min_{x \in \mathbb{R}^d} f(x), \qquad f(x) = \mathbb{E}_{z}[F(x;z)],
\]
where $f$ is smooth and (possibly) non-convex.

\paragraph{State variables.}
\begin{itemize}
  \item $x_t \in \mathbb{R}^d$: iterate at step $t$.
  \item $v_t \in \mathbb{R}^d$: recursive gradient/variance estimator.
  \item $\hat v_t \in \mathbb{R}^d$: bias-corrected estimator.
\end{itemize}

\paragraph{Hyperparameters.}
\begin{itemize}
  \item $\alpha \in (0,1]$: smoothing/recursion parameter.
  \item $\eta > 0$: base step size.
  \item $\epsilon > 0$: small constant ensuring numerical stability.
  \item $T \in \mathbb{N}$: number of iterations.
\end{itemize}

\paragraph{Update rules.}
\begin{align}
v_t        &= (1-\alpha)\, v_{t-1} + \alpha\, g_t,
   \qquad g_t = \nabla f(x_t;z_t), \label{eq:vt}\\
\hat v_t   &= \frac{v_t}{1-(1-\alpha)^t}, \label{eq:bias}\\
\eta_t     &= \frac{\eta}{\sqrt{\|\hat v_t\|^2 + \epsilon}}, \label{eq:step}\\
x_{t+1}    &= x_t - \eta_t\, \hat v_t. \label{eq:upd}
\end{align}

Here $g_t$ is an unbiased stochastic gradient: $\mathbb{E}[g_t \mid x_t] = \nabla f(x_t)$.
The norm-adaptive step \eqref{eq:step} normalizes the update so that
$\|x_{t+1}-x_t\| = \eta_t\|\hat v_t\| = \eta \|\hat v_t\|/\sqrt{\|\hat v_t\|^2+\epsilon} \le \eta$.

%==========================================================
\section*{Pseudocode}

\begin{algorithm}[H]
\caption{RVE-NA: Recursive Variance Estimator with Norm-Adaptive Step}
\begin{algorithmic}[1]
\State \textbf{Input:} $x_1$, $\alpha\in(0,1]$, $\eta>0$, $\epsilon>0$, horizon $T$
\State \textbf{Initialize:} $v_0 \gets 0$
\For{$t = 1, 2, \dots, T$}
    \State Sample $z_t$; compute $g_t \gets \nabla f(x_t;z_t)$
    \State $v_t \gets (1-\alpha) v_{t-1} + \alpha g_t$ \Comment{recursive estimator}
    \State $\hat v_t \gets v_t / \big(1-(1-\alpha)^t\big)$ \Comment{bias correction}
    \State $\eta_t \gets \eta / \sqrt{\|\hat v_t\|^2 + \epsilon}$ \Comment{norm-adaptive step}
    \State $x_{t+1} \gets x_t - \eta_t \hat v_t$
\EndFor
\State \textbf{Termination:} stop at $T$ iterations or when $\|\hat v_t\| \le \mathrm{tol}$
\State \textbf{Output:} $x_R$ where $R$ is chosen uniformly at random from $\{1,\dots,T\}$
\end{algorithmic}
\end{algorithm}

%==========================================================
\section*{Dry Run Example}

Let $f(w) = (w-3)^2$, so $\nabla f(w) = 2(w-3)$.
Take $w_0 = 0$ (i.e. $x_1 = 0$), $\eta = 0.1$, $\alpha = 0.5$, $\epsilon = 10^{-8}$.
We treat gradients as exact (deterministic) for illustration. Scalar case: $\|\cdot\|=|\cdot|$.

\paragraph{Iteration $t=1$.}
\begin{align*}
g_1 &= 2(0-3) = -6,\\
v_1 &= (1-0.5)\cdot 0 + 0.5\cdot(-6) = -3,\\
\hat v_1 &= \frac{-3}{1-(0.5)^1} = \frac{-3}{0.5} = -6,\\
\eta_1 &= \frac{0.1}{\sqrt{(-6)^2+10^{-8}}} = \frac{0.1}{6.0000\ldots} \approx 0.016667,\\
x_2 &= 0 - 0.016667\cdot(-6) = 0 + 0.1 = 0.1,\\
f(x_2) &= (0.1-3)^2 = 8.41.
\end{align*}

\paragraph{Iteration $t=2$.}
\begin{align*}
g_2 &= 2(0.1-3) = -5.8,\\
v_2 &= 0.5\cdot(-3) + 0.5\cdot(-5.8) = -4.4,\\
\hat v_2 &= \frac{-4.4}{1-(0.5)^2} = \frac{-4.4}{0.75} \approx -5.8667,\\
\eta_2 &= \frac{0.1}{\sqrt{(-5.8667)^2+10^{-8}}} \approx \frac{0.1}{5.8667} \approx 0.017045,\\
x_3 &= 0.1 - 0.017045\cdot(-5.8667) = 0.1 + 0.1 = 0.2,\\
f(x_3) &= (0.2-3)^2 = 7.84.
\end{align*}

\paragraph{Iteration $t=3$.}
\begin{align*}
g_3 &= 2(0.2-3) = -5.6,\\
v_3 &= 0.5\cdot(-4.4) + 0.5\cdot(-5.6) = -5.0,\\
\hat v_3 &= \frac{-5.0}{1-(0.5)^3} = \frac{-5.0}{0.875} \approx -5.7143,\\
\eta_3 &= \frac{0.1}{\sqrt{(-5.7143)^2+10^{-8}}} \approx \frac{0.1}{5.7143} \approx 0.0175,\\
x_4 &= 0.2 - 0.0175\cdot(-5.7143) = 0.2 + 0.1 = 0.3,\\
f(x_4) &= (0.3-3)^2 = 7.29.
\end{align*}

\noindent
The objective decreases monotonically ($8.41 \to 7.84 \to 7.29$) and the iterate moves
toward the minimizer $w^\star=3$. The norm-adaptive step produces a fixed displacement of
magnitude $\approx \eta = 0.1$ per step (since $\epsilon$ is negligible), illustrating the
normalized-update behavior.

%==========================================================
\section*{Assumptions}

\begin{assumption}[Lower boundedness]\label{as:lb}
$f$ is bounded below: $f^\star := \inf_x f(x) > -\infty$.
\end{assumption}

\begin{assumption}[$L$-smoothness]\label{as:smooth}
$f$ is differentiable and $\nabla f$ is $L$-Lipschitz:
\[
\|\nabla f(x) - \nabla f(y)\| \le L\|x-y\| \quad \forall x,y.
\]
Equivalently, the descent lemma holds:
\[
f(y) \le f(x) + \langle \nabla f(x), y-x\rangle + \tfrac{L}{2}\|y-x\|^2.
\]
\end{assumption}

\begin{assumption}[Unbiased gradients]\label{as:unbiased}
$\mathbb{E}[g_t \mid x_t] = \nabla f(x_t)$.
\end{assumption}

\begin{assumption}[Bounded estimator error / variance]\label{as:var}
The bias-corrected recursive estimator tracks the true gradient with bounded
mean-squared error: there exists $\sigma^2 \ge 0$ such that for all $t$,
\[
\mathbb{E}\big[\|\hat v_t - \nabla f(x_t)\|^2\big] \le \sigma_t^2,
\qquad
\frac{1}{T}\sum_{t=1}^T \sigma_t^2 \le \sigma^2 .
\]
For SARAH/SPIDER-type recursions with mini-batch size $b$ and inner-loop length $m$,
$\sigma_t^2 = O\!\big(\tfrac{\alpha}{b} + \tfrac{(1-\alpha) L^2}{b}\sum_{s\le t}\|x_s-x_{s-1}\|^2\big)$,
which is controlled by the normalized steps.
\end{assumption}

\begin{assumption}[Bounded gradients]\label{as:bg}
$\|\nabla f(x)\| \le G$ and $\|\hat v_t\| \le \hat G$ for all $x,t$.
\end{assumption}

%==========================================================
\section*{Main Convergence Theorem}

\begin{theorem}[Non-convex convergence rate]\label{thm:main}
Under Assumptions~\ref{as:lb}--\ref{as:bg}, run RVE-NA for $T$ iterations with
constant base step $\eta>0$ satisfying $\eta \le \frac{\sqrt{\epsilon}}{2L}$ (in the regime
$\|\hat v_t\|^2 \le \epsilon$ effective bound, see Step 9). Let $R$ be sampled uniformly
from $\{1,\dots,T\}$. Then
\[
\mathbb{E}\big[\|\nabla f(x_R)\|^2\big]
\;\le\;
\frac{2\sqrt{\hat G^2+\epsilon}\,\big(f(x_1)-f^\star\big)}{\eta\,T}
\;+\;
\frac{L\eta\, \hat G^2}{\sqrt{\epsilon}}
\;+\;
2\sigma^2 .
\]
In particular, with $\eta = \Theta(1/\sqrt{T})$ and $\sigma^2 = O(1/\sqrt{T})$ (achievable
by SARAH/SPIDER variance control), 
\[
\frac{1}{T}\sum_{t=1}^T \mathbb{E}\big[\|\nabla f(x_t)\|^2\big]
\;=\; O\!\left(\frac{1}{\sqrt{T}}\right).
\]
\end{theorem}

%==========================================================
\section*{Theoretical Properties}

\begin{itemize}
\item \textbf{Stationarity guarantee.} The bound controls the average squared gradient norm,
the standard non-convex stationarity measure.
\item \textbf{Stability.} The norm-adaptive step \eqref{eq:step} guarantees bounded
displacement $\|x_{t+1}-x_t\| \le \eta$, preventing divergence even with large gradients.
\item \textbf{Hyperparameter sensitivity.} Smaller $\alpha$ yields stronger smoothing
(lower variance, more bias); bias correction \eqref{eq:bias} removes the initialization bias
$(1-\alpha)^t$. The constant $\epsilon$ trades off adaptivity vs.\ stability.
\item \textbf{Comparison to baselines.} Compared with SGD ($O(1/\sqrt{T})$ for stationarity
with $\eta=\Theta(1/\sqrt{T})$), the recursive estimator reduces $\sigma^2$, and SARAH/SPIDER
variance control can push the complexity toward $O(1/T)$-type behavior when $\sigma_t^2$ is
driven down by the bounded normalized steps.
\end{itemize}

%==========================================================
\section*{Discussion}
The algorithm combines three ingredients: (i) a recursive (momentum-like) gradient estimator
that lowers variance, (ii) Adam-style bias correction that removes transient initialization
error, and (iii) a normalized step that ensures stable, bounded updates in the non-convex
landscape. The proof below isolates the bias of $\hat v_t$ relative to $\nabla f(x_t)$ and
shows that, provided the estimator variance is summable/controlled, the method reaches an
approximate stationary point.

%==========================================================
\section*{Preliminaries (Definitions and Lemmas)}

\begin{lemma}[Descent Lemma]\label{lem:descent}
Under Assumption~\ref{as:smooth}, for any $x,y$,
\[
f(y) \le f(x) + \langle \nabla f(x), y-x\rangle + \frac{L}{2}\|y-x\|^2.
\]
\end{lemma}

\begin{lemma}[Step-size bounds]\label{lem:step}
Under Assumption~\ref{as:bg}, the adaptive step \eqref{eq:step} satisfies
\[
\eta_t = \frac{\eta}{\sqrt{\|\hat v_t\|^2+\epsilon}}
\in \left[\frac{\eta}{\sqrt{\hat G^2+\epsilon}},\ \frac{\eta}{\sqrt{\epsilon}}\right].
\]
\end{lemma}
\begin{proof}
Since $0 \le \|\hat v_t\|^2 \le \hat G^2$ (Assumption~\ref{as:bg}), we have
$\sqrt{\epsilon} \le \sqrt{\|\hat v_t\|^2+\epsilon} \le \sqrt{\hat G^2+\epsilon}$;
taking reciprocals and multiplying by $\eta$ gives the claim.
\end{proof}

\begin{lemma}[Polarization / Young's inequality]\label{lem:young}
For any vectors $a,b$ and $\rho>0$:
$\langle a,b\rangle \ge -\tfrac{1}{2}\|a\|^2 - \tfrac{1}{2}\|b\|^2$, and more generally
$2\langle a,b\rangle \le \rho\|a\|^2 + \rho^{-1}\|b\|^2$. Also
$\|a+b\|^2 \le 2\|a\|^2 + 2\|b\|^2$.
\end{lemma}

%==========================================================
\section*{Main Proof (Step-by-Step Derivation)}

\noindent\textbf{[Step 1] Apply the Descent Lemma.}
By Lemma~\ref{lem:descent} with $x=x_t$, $y=x_{t+1}=x_t-\eta_t\hat v_t$:
\[
f(x_{t+1}) \le f(x_t) - \eta_t \langle \nabla f(x_t), \hat v_t\rangle
+ \frac{L}{2}\eta_t^2 \|\hat v_t\|^2.
\tag{S1}
\]

\noindent\textbf{[Step 2] Decompose the inner product around the true gradient.}
Write $\hat v_t = \nabla f(x_t) + e_t$ where $e_t := \hat v_t - \nabla f(x_t)$ is the
estimator error. Then
\[
\langle \nabla f(x_t), \hat v_t\rangle
= \|\nabla f(x_t)\|^2 + \langle \nabla f(x_t), e_t\rangle.
\tag{S2}
\]

\noindent\textbf{[Step 3] Substitute into (S1).}
\[
f(x_{t+1}) \le f(x_t) - \eta_t \|\nabla f(x_t)\|^2
- \eta_t \langle \nabla f(x_t), e_t\rangle
+ \frac{L}{2}\eta_t^2 \|\hat v_t\|^2.
\tag{S3}
\]

\noindent\textbf{[Step 4] Bound the cross term via Young's inequality.}
Using Lemma~\ref{lem:young} with $\rho=1$:
\[
-\langle \nabla f(x_t), e_t\rangle
\le \tfrac12 \|\nabla f(x_t)\|^2 + \tfrac12 \|e_t\|^2.
\tag{S4}
\]
Hence
\[
-\eta_t\langle \nabla f(x_t), e_t\rangle
\le \frac{\eta_t}{2}\|\nabla f(x_t)\|^2 + \frac{\eta_t}{2}\|e_t\|^2.
\]

\noindent\textbf{[Step 5] Combine (S3)–(S4).}
\[
f(x_{t+1}) \le f(x_t) - \frac{\eta_t}{2}\|\nabla f(x_t)\|^2
+ \frac{\eta_t}{2}\|e_t\|^2
+ \frac{L}{2}\eta_t^2 \|\hat v_t\|^2.
\tag{S5}
\]

\noindent\textbf{[Step 6] Use step-size lower bound on the negative term.}
By Lemma~\ref{lem:step}, $\eta_t \ge \eta/\sqrt{\hat G^2+\epsilon} =: \eta_{\min}$.
Since the coefficient $-\tfrac{\eta_t}{2}\|\nabla f(x_t)\|^2 \le -\tfrac{\eta_{\min}}{2}\|\nabla f(x_t)\|^2$,
\[
f(x_{t+1}) \le f(x_t) - \frac{\eta_{\min}}{2}\|\nabla f(x_t)\|^2
+ \frac{\eta_t}{2}\|e_t\|^2
+ \frac{L}{2}\eta_t^2 \|\hat v_t\|^2.
\tag{S6}
\]

\noindent\textbf{[Step 7] Use step-size upper bound on the error/curvature terms.}
By Lemma~\ref{lem:step}, $\eta_t \le \eta/\sqrt{\epsilon} =: \eta_{\max}$,
and by Assumption~\ref{as:bg} $\|\hat v_t\|^2 \le \hat G^2$. Therefore
\[
\frac{\eta_t}{2}\|e_t\|^2 \le \frac{\eta_{\max}}{2}\|e_t\|^2,
\qquad
\frac{L}{2}\eta_t^2 \|\hat v_t\|^2 \le \frac{L}{2}\eta_{\max}^2 \hat G^2
= \frac{L\eta^2 \hat G^2}{2\epsilon}.
\tag{S7}
\]
Substituting:
\[
f(x_{t+1}) \le f(x_t) - \frac{\eta_{\min}}{2}\|\nabla f(x_t)\|^2
+ \frac{\eta_{\max}}{2}\|e_t\|^2
+ \frac{L\eta^2 \hat G^2}{2\epsilon}.
\tag{S7'}
\]

\noindent\textbf{[Step 8] Take expectations.}
Taking total expectation and using Assumption~\ref{as:var}
($\mathbb{E}\|e_t\|^2 \le \sigma_t^2$):
\[
\mathbb{E}[f(x_{t+1})] \le \mathbb{E}[f(x_t)]
- \frac{\eta_{\min}}{2}\mathbb{E}\|\nabla f(x_t)\|^2
+ \frac{\eta_{\max}}{2}\sigma_t^2
+ \frac{L\eta^2 \hat G^2}{2\epsilon}.
\tag{S8}
\]

\noindent\textbf{[Step 9] Rearrange to isolate the gradient term.}
\[
\frac{\eta_{\min}}{2}\mathbb{E}\|\nabla f(x_t)\|^2
\le \mathbb{E}[f(x_t)] - \mathbb{E}[f(x_{t+1})]
+ \frac{\eta_{\max}}{2}\sigma_t^2
+ \frac{L\eta^2 \hat G^2}{2\epsilon}.
\tag{S9}
\]

\noindent\textbf{[Step 10] Sum over $t=1,\dots,T$ (telescoping).}
\[
\frac{\eta_{\min}}{2}\sum_{t=1}^T \mathbb{E}\|\nabla f(x_t)\|^2
\le \big(f(x_1)-\mathbb{E}[f(x_{T+1})]\big)
+ \frac{\eta_{\max}}{2}\sum_{t=1}^T \sigma_t^2
+ \frac{L\eta^2 \hat G^2}{2\epsilon}T.
\tag{S10}
\]
The telescoping uses
$\sum_{t=1}^T \big(\mathbb{E}[f(x_t)]-\mathbb{E}[f(x_{t+1})]\big)
= f(x_1)-\mathbb{E}[f(x_{T+1})]$.

\noindent\textbf{[Step 11] Apply lower boundedness.}
Since $\mathbb{E}[f(x_{T+1})] \ge f^\star$ (Assumption~\ref{as:lb}),
$f(x_1)-\mathbb{E}[f(x_{T+1})] \le f(x_1)-f^\star$. Thus
\[
\frac{\eta_{\min}}{2}\sum_{t=1}^T \mathbb{E}\|\nabla f(x_t)\|^2
\le \big(f(x_1)-f^\star\big)
+ \frac{\eta_{\max}}{2}\sum_{t=1}^T \sigma_t^2
+ \frac{L\eta^2 \hat G^2}{2\epsilon}T.
\tag{S11}
\]

\noindent\textbf{[Step 12] Divide by $\tfrac{\eta_{\min}}{2}T$.}
\[
\frac{1}{T}\sum_{t=1}^T \mathbb{E}\|\nabla f(x_t)\|^2
\le \frac{2\big(f(x_1)-f^\star\big)}{\eta_{\min} T}
+ \frac{\eta_{\max}}{\eta_{\min}}\cdot\frac{1}{T}\sum_{t=1}^T \sigma_t^2
+ \frac{L\eta^2 \hat G^2}{\epsilon\,\eta_{\min}}.
\tag{S12}
\]

\noindent\textbf{[Step 13] Insert random-index reformulation.}
Since $R\sim \mathrm{Unif}\{1,\dots,T\}$,
$\mathbb{E}\|\nabla f(x_R)\|^2 = \frac{1}{T}\sum_{t=1}^T \mathbb{E}\|\nabla f(x_t)\|^2$,
which equals the left side of (S12).

%==========================================================
\section*{Rate Derivation (With Constants)}

\noindent\textbf{[Step 14] Substitute explicit step-size constants.}
Recall $\eta_{\min} = \eta/\sqrt{\hat G^2+\epsilon}$ and $\eta_{\max}=\eta/\sqrt{\epsilon}$, so
\[
\frac{1}{\eta_{\min}} = \frac{\sqrt{\hat G^2+\epsilon}}{\eta},
\qquad
\frac{\eta_{\max}}{\eta_{\min}} = \frac{\sqrt{\hat G^2+\epsilon}}{\sqrt{\epsilon}},
\qquad
\frac{\eta^2}{\eta_{\min}} = \eta\sqrt{\hat G^2+\epsilon}.
\]
Plugging into (S12):
\begin{align}
\mathbb{E}\|\nabla f(x_R)\|^2
&\le \frac{2\sqrt{\hat G^2+\epsilon}\,(f(x_1)-f^\star)}{\eta\,T}
+ \frac{\sqrt{\hat G^2+\epsilon}}{\sqrt{\epsilon}}\cdot\frac{1}{T}\sum_{t=1}^T\sigma_t^2
\notag\\
&\quad + \frac{L\,\hat G^2 \sqrt{\hat G^2+\epsilon}}{\epsilon}\,\eta.
\tag{S14}
\end{align}

\noindent\textbf{[Step 15] Simplify using $\tfrac1T\sum\sigma_t^2 \le \sigma^2$ and constant absorption.}
Writing $C_0 = \sqrt{\hat G^2+\epsilon}$ and noting $C_0/\sqrt{\epsilon}\ge 1$, a slightly
loosened but cleaner bound (matching Theorem~\ref{thm:main}, up to the constant $C_0$) is
\[
\mathbb{E}\|\nabla f(x_R)\|^2
\le \frac{2 C_0\,(f(x_1)-f^\star)}{\eta\,T}
+ \frac{C_0}{\sqrt{\epsilon}}\,\sigma^2
+ \frac{L C_0\,\hat G^2}{\epsilon}\,\eta.
\tag{S15}
\]

\noindent\textbf{[Step 16] Optimize the base step $\eta$.}
The $\eta$-dependent terms are $A/\eta + B\eta$ with
$A = 2C_0(f(x_1)-f^\star)/T$ and $B = LC_0\hat G^2/\epsilon$.
Minimizing over $\eta$ gives $\eta^\star=\sqrt{A/B}$ and minimal value $2\sqrt{AB}$:
\[
2\sqrt{AB}
= 2\sqrt{\frac{2C_0(f(x_1)-f^\star)}{T}\cdot \frac{LC_0\hat G^2}{\epsilon}}
= 2C_0\hat G\sqrt{\frac{2L(f(x_1)-f^\star)}{\epsilon\,T}}
= O\!\left(\frac{1}{\sqrt T}\right).
\tag{S16}
\]
Therefore, with the optimal (or $\eta=\Theta(1/\sqrt T)$) choice,
\[
\mathbb{E}\|\nabla f(x_R)\|^2
\le O\!\left(\frac{1}{\sqrt T}\right) + \frac{C_0}{\sqrt\epsilon}\sigma^2 .
\tag{S17}
\]

\noindent\textbf{[Step 17] Variance-reduction (SARAH/SPIDER) regime.}
If the recursive estimator achieves $\tfrac1T\sum_t\sigma_t^2 = O(1/\sqrt T)$ (as in
SARAH/SPIDER with batch/length tuned to the bounded normalized steps,
Assumption~\ref{as:var}), then both terms in (S17) are $O(1/\sqrt T)$, yielding
\[
\boxed{\;\mathbb{E}\big[\|\nabla f(x_R)\|^2\big] = O\!\left(\frac{1}{\sqrt{T}}\right).\;}
\]

%==========================================================
\section*{Special Cases}

\paragraph{Case $\alpha=1$ (plain normalized SGD).}
Then $v_t=g_t$, $\hat v_t = g_t/(1-0)=g_t$, and the method reduces to normalized stochastic
gradient descent with adaptive step $\eta_t=\eta/\sqrt{\|g_t\|^2+\epsilon}$. The bound (S15)
holds with $\sigma_t^2 = \mathbb{E}\|g_t-\nabla f(x_t)\|^2$ (raw variance), recovering the
standard non-convex SGD rate $O(1/\sqrt T)$ plus a non-vanishing variance floor.

\paragraph{Case $\sigma_t^2 \to 0$ (exact/full gradients).}
If $g_t=\nabla f(x_t)$ (deterministic), then $e_t \to 0$ after bias correction and the
variance term vanishes. The bound becomes
\[
\frac{1}{T}\sum_{t=1}^T \|\nabla f(x_t)\|^2
\le \frac{2C_0(f(x_1)-f^\star)}{\eta T} + \frac{L C_0 \hat G^2}{\epsilon}\eta,
\]
giving a clean $O(1/T)$ in the leading term for fixed small $\eta$ (gradient-norm
stationarity), consistent with deterministic normalized gradient descent.

\paragraph{Case small $\epsilon$ (strong normalization).}
As $\epsilon\to 0$, $\eta_t \approx \eta/\|\hat v_t\|$ enforces unit-scaled steps
($\|x_{t+1}-x_t\|\approx\eta$), maximizing stability but inflating the curvature constant
$B=LC_0\hat G^2/\epsilon$; the step $\eta$ must then be reduced accordingly, illustrating
the stability/accuracy trade-off encoded by $\epsilon$.

\end{document}